% Generated by task tex-boundary from dossiers/boundary-adaptor.md, appendix (the eight probes);
% sources and outputs are copied verbatim from the dossier, which reproduces the files under
% .claude/reports/blueprint-review/probes/boundary-adaptor/ (checked identical).
% This fragment is read inside \section{The boundary with leanISA} of sections/b-probes.tex, so
% each probe is a \subsection.
% Breakable inline code for this fragment's tables and prose (LuaLaTeX): \bndc{...} sets its
% argument verbatim in the monospace font, with break points after punctuation and between the
% words of a camel-case identifier. Defined once, whichever boundary fragment is read first.
\ifdefined\bndc\else
\directlua{
  bndbreak = function(s)
    local bs = string.char(92)
    local hash = string.char(35)
    s = s:gsub(hash .. hash, hash)
    local after = { ["."] = true, ["_"] = true, ["/"] = true, [":"] = true, [","] = true,
      ["-"] = true, [")"] = true, ["]"] = true, ["="] = true, ["|"] = true }
    local prev = ""
    for _, c in utf8.codes(s) do
      local ch = utf8.char(c)
      if ch:match("^[A-Z]$") and prev:match("^[a-z]$") then
        tex.sprint(bs .. "penalty400{}")
      end
      tex.sprint(-2, ch)
      if after[ch] or c > 127 then
        tex.sprint(bs .. "allowbreak{}")
      end
      prev = ch
    end
  end
}
\newcommand{\bndc}[1]{\texttt{\directlua{bndbreak("\luaescapestring{\detokenize{#1}}")}}}
\fi

\label{sec:gen-boundary-probes}
\begin{sloppypar}
The eight probes of the boundary dossier. Every file is under
\bndc{.claude/reports/blueprint-review/probes/boundary-adaptor/} and was run from the repository
root under the shared lock; the recorded output follows each source, with the exit status on its
last line. An empty output with \code{exit=0} means every \bndc{#guard} passed and every
\code{example} elaborated. \code{Shapes}, \code{KnownColumn}, \code{Transport},
\code{UniverseFail}, \code{DefInstance} and \code{PolyBridge} ran on 2026-09-29 against
\code{main} at \code{b435631} with the old pins (ArkLib \code{dca90385}, CompPoly \code{3468b38c},
Clean \code{93c9d1ef}). Two of them use numerals of \code{K} other than \code{0} and \code{1}
(\code{KnownColumn}: \code{2}; \code{PolyBridge}: \code{3, 5, 7}), which at the new CompPoly pin
mean \code{0}, \code{1}, \code{1}, \code{1}; their conclusions stand at the old pin, and both were
re-run at the new pins, with \code{K.ofBits n} in place of the numerals, as \code{KnownColumnNew}
and \code{PolyBridgeNew} on 2026-09-30. The other four were not re-run; nothing they rest on
changed between the two revisions of the proof-system sources, which differ in proofs and
imports only. Line numbers are shown beside sources and outputs.
\end{sloppypar}

\subsection{\texorpdfstring{\protect\code{Shapes}}{Shapes}}\label{sec:gen-boundary-probe-shapes}

\paragraph{What it establishes.} The two relations, their universes and the public input. It
prints the types of \bndc{EnsembleWitness}, \bndc{Air.Flat.Table} and \bndc{Component}, all in
\bndc{Type 1} (output lines 1--4), of \bndc{M3Rel} and \bndc{Column}, in \bndc{Type}, of
\bndc{Refinement}, whose two witness universes are independent (output lines 10--13), and of
\bndc{Extractor.Straightline.map}, whose \bndc{WitIn'} is in \bndc{Type}; the statements of
\bndc{SatisfiedBy} (output lines 29--55), \bndc{Caps}, \bndc{PublicInput} and its two words
(69--78), \bndc{ValidExecution}, \bndc{HasPublicBoundary}, \bndc{Program}, \bndc{M3Holds} and
\bndc{M3Instance} (100--134), \bndc{PublicLine}; and what \bndc{w.Constraints} unfolds to, Clean's
\bndc{ConstraintsHold} with constraints and lookups only, no bus guarantees (154--159). Its four
\bndc{example}s prove that the top limb of both public words is zero for every input (source lines
57--61), that a word of the image whose top limb is nonzero is not \bndc{p.word0}, whatever the
input (63--69), and that the three public lines give \bndc{word0_eq} (71--76); the dossier cites
these source lines as 45--63, twelve lines too low. The axioms of every arithmetization
declaration the chain will use, of the two master theorems and of \bndc{bytecodeColumn_eval},
\bndc{idxColumn_eval} are the kernel's three; \bndc{Refinement.map_option_valid} depends on two
(output lines 168--187). Used in \cref{sec:gen-boundary-objects,sec:gen-boundary-toplimb}.

\paragraph{Revision.} Run on 2026-09-29 against leanerVM \code{main} at \code{b435631} with the old pins (ArkLib \code{dca90385}, CompPoly \code{3468b38c}, Clean \code{93c9d1ef}). Exit 0. Not re-run at the new pins; nothing it rests on changed between the two revisions of the proof-system sources (proofs and imports only).

\paragraph{Source.}
\begin{leancode}[numbers=left,numbersep=3pt]
/-
  Probe (boundary-adaptor): the two relations, their universes, and the public input.
  Scratch work for the blueprint review; not part of the repository.
  Run from the repository root:
    flock .claude/reports/blueprint-review/logs/lean.lock \
      lake env lean .claude/reports/blueprint-review/probes/boundary-adaptor/Shapes.lean
-/
import LeanerVM

open LeanerVM LeanerVM.Parameters LeanerVM.Semantics LeanerVM.Arithmetization
open LeanerVM.Protocol
open Air.Flat (Ensemble EnsembleWitness Component)

set_option pp.universes true

/-! ## 1. Universes: the witness of `SatisfiedBy` against the witness of `M3Holds` -/

#check @EnsembleWitness
#check @Air.Flat.Table
#check @Air.Flat.Component
#check @leanIsaEnsemble
#check @SatisfiedBy
#check @M3Holds
#check @M3Rel
#check @Column
#check @Refinement
#check @Refinement.map_option_valid
#check @Extractor.Straightline.map

set_option pp.universes false

/-! ## 2. The statement types as written -/

#print SatisfiedBy
#print Caps
#print PublicInput
#print PublicInput.word0
#print PublicInput.word1
#print PublicIO
#print ValidExecution
#print HasPublicBoundary
#print Program
#print M3Holds
#print M3Instance
#print PublicLine

/-! ## 3. Which `ConstraintsHold` does `w.Constraints` unfold to? -/

#print Air.Flat.EnsembleWitness.Constraints
#print Air.Flat.Table.Constraints
#print Operations.ConstraintsHold
#print Air.Flat.EnsembleWitness.interactions
#print ProverData

/-! ## 4. The top limb of a public word is zero by construction -/

example (p : PublicInput) : p.word0.limb 2 = 0 := by
  simp [PublicInput.word0]

example (p : PublicInput) : p.word1.limb 2 = 0 := by
  simp [PublicInput.word1]

/-- A word of the image whose top limb is nonzero is not a public word, whatever the input: the
fact `word0_eq` needs from the third public line. -/
example (p : PublicInput) (a b c : K) (hc : c ≠ 0) : E.ofLimbs a b c ≠ p.word0 := by
  intro h
  have := congrArg (fun z : E ↦ z.limb 2) h
  simp [PublicInput.word0] at this
  exact hc this

/-- With the top limb zero and the two low limbs those of the statement, the word is the public
word: the three public lines give `word0_eq`. -/
example (p : PublicInput) (a b c : K) (ha : a = p.word0.limb 0) (hb : b = p.word0.limb 1)
    (hc : c = 0) : E.ofLimbs a b c = p.word0 := by
  subst ha hb hc
  simp [PublicInput.word0]

/-! ## 5. Axioms of what the chain will use on the arithmetization side, and of the master
theorems -/

#print axioms assumptions_of_blake2sRowsValid
#print axioms memRowOf_bindings
#print axioms bytecodeRowOf_bindings
#print axioms bytecodeRowOf_decodes
#print axioms verifier_pull_eval
#print axioms verifier_push_eval
#print axioms rowAt_toElements
#print axioms decode_eq_some_iff
#print axioms assignmentRepresents_image
#print axioms xorTable
#print axioms jumpTable
#print axioms blake2sTable
#print axioms memTable
#print axioms bytecodeTable
#print axioms leanIsaVerifier
#print axioms piop_rbrKnowledgeSoundness
#print axioms piop_perfectCompleteness
#print axioms Refinement.map_option_valid
#print axioms bytecodeColumn_eval
#print axioms idxColumn_eval
\end{leancode}
\src{.claude/reports/blueprint-review/probes/boundary-adaptor/Shapes.lean}

\paragraph{Command} (from the repository root).
\begin{shellcode}
flock .claude/reports/blueprint-review/logs/lean.lock \
  lake env lean .claude/reports/blueprint-review/probes/boundary-adaptor/Shapes.lean
\end{shellcode}

\paragraph{Recorded output} (\file{Shapes.out}).
\begin{srccode}[numbers=left,numbersep=3pt]
@EnsembleWitness : {F : Type} →
  [inst : FiniteField F] → {PublicIO : TypeMap} → [inst_1 : ProvableType PublicIO] → Ensemble F PublicIO → Type 1
Air.Flat.Table : (F : Type) → [FiniteField F] → Type 1
Component : (F : Type) → [FiniteField F] → Type 1
leanIsaEnsemble : Program → Ensemble K PublicIO
SatisfiedBy : (prog : Program) → PublicInput → EnsembleWitness (leanIsaEnsemble prog) → Prop
M3Holds : (I : M3Instance) → I.Stmt → Column I.μ → Prop
M3Rel : (I : M3Instance) → Set.{0} (Prod.{0, 0} (Prod.{0, 0} I.Stmt ((i : Fin 0) → NoOracle i)) (Column I.μ))
Column : ℕ → Type
@Refinement.{u_1,
    u_2} : {Stmt : Type} →
  {W₁ : Type u_1} →
    {W₂ : Type u_2} → Set.{u_1} (Prod.{0, u_1} Stmt W₁) → Set.{u_2} (Prod.{0, u_2} Stmt W₂) → Type (max u_1 u_2)
@Refinement.map_option_valid.{u_1,
    u_2} : ∀ {Stmt : Type} {W₁ : Type u_1} {W₂ : Type u_2} {R : Set.{u_1} (Prod.{0, u_1} Stmt W₁)}
  {S : Set.{u_2} (Prod.{0, u_2} Stmt W₂)} (f : Refinement.{u_1, u_2} R S) (x : Stmt) (w? : Option.{u_1} W₁),
  (∀ (w : W₁), Membership.mem.{u_1, u_1} w? w → Membership.mem.{u_1, u_1} R (Prod.mk.{0, u_1} x w)) →
    ∀ (w' : W₂),
      Membership.mem.{u_2, u_2} (Option.map.{u_1, u_2} (Refinement.map.{u_1, u_2} f x) w?) w' →
        Membership.mem.{u_2, u_2} S (Prod.mk.{0, u_2} x w')
@Extractor.Straightline.map.{u_1} : {ι : Type} →
  {oSpec : OracleSpec.{0, u_1} ι} →
    {StmtIn WitIn WitIn' WitOut : Type} →
      {n : ℕ} →
        {pSpec : ProtocolSpec n} →
          (StmtIn → WitIn → WitIn') →
            Extractor.Straightline.{u_1} oSpec StmtIn WitIn WitOut pSpec →
              Extractor.Straightline.{u_1} oSpec StmtIn WitIn' WitOut pSpec
structure LeanerVM.Arithmetization.SatisfiedBy (prog : Program) (input : PublicInput)
  (w : EnsembleWitness (leanIsaEnsemble prog)) : Prop
number of parameters: 3
fields:
  LeanerVM.Arithmetization.SatisfiedBy.public_input_eq : w.publicInput = PublicIO.ofInput input
  LeanerVM.Arithmetization.SatisfiedBy.constraints : w.Constraints
  LeanerVM.Arithmetization.SatisfiedBy.state_balanced : BalancedPair w StatePull.toRaw StatePush.toRaw
  LeanerVM.Arithmetization.SatisfiedBy.mem_balanced : BalancedPair w MemPull.toRaw MemPush.toRaw
  LeanerVM.Arithmetization.SatisfiedBy.bytecode_balanced : BalancedPair w BytecodePull.toRaw BytecodePush.toRaw
  LeanerVM.Arithmetization.SatisfiedBy.counts_nonzero : CountsNonzero w
  LeanerVM.Arithmetization.SatisfiedBy.caps : Caps w
  LeanerVM.Arithmetization.SatisfiedBy.index_columns : IndexColumnsAreRowIndices w
  LeanerVM.Arithmetization.SatisfiedBy.seed_rows : SeedRowsAreTheImage w
  LeanerVM.Arithmetization.SatisfiedBy.bytecode_rows : BytecodeRowsAreTheProgram prog w
  LeanerVM.Arithmetization.SatisfiedBy.blake2s_valid : Blake2sRowsValid w
  LeanerVM.Arithmetization.SatisfiedBy.word0_eq : (imageOf w.data).snd.read (gpow 0) = some input.word0
  LeanerVM.Arithmetization.SatisfiedBy.word1_eq : (imageOf w.data).snd.read (gpow 1) = some input.word1
constructor:
  LeanerVM.Arithmetization.SatisfiedBy.mk {prog : Program} {input : PublicInput}
    {w : EnsembleWitness (leanIsaEnsemble prog)} (public_input_eq : w.publicInput = PublicIO.ofInput input)
    (constraints : w.Constraints) (state_balanced : BalancedPair w StatePull.toRaw StatePush.toRaw)
    (mem_balanced : BalancedPair w MemPull.toRaw MemPush.toRaw)
    (bytecode_balanced : BalancedPair w BytecodePull.toRaw BytecodePush.toRaw) (counts_nonzero : CountsNonzero w)
    (caps : Caps w) (index_columns : IndexColumnsAreRowIndices w) (seed_rows : SeedRowsAreTheImage w)
    (bytecode_rows : BytecodeRowsAreTheProgram prog w) (blake2s_valid : Blake2sRowsValid w)
    (word0_eq : (imageOf w.data).snd.read (gpow 0) = some input.word0)
    (word1_eq : (imageOf w.data).snd.read (gpow 1) = some input.word1) : SatisfiedBy prog input w
structure LeanerVM.Arithmetization.Caps {prog : Program} (w : EnsembleWitness (leanIsaEnsemble prog)) : Prop
number of parameters: 2
fields:
  LeanerVM.Arithmetization.Caps.minLogMem_le : minLogMem ≤ (imageOf w.data).fst
  LeanerVM.Arithmetization.Caps.le_maxLogMem : (imageOf w.data).fst ≤ maxLogMem
  LeanerVM.Arithmetization.Caps.heights : ∀ t ∈ w.tables, ∃ τ ≤ maxLogRows, t.table.length = 2 ^ τ
  LeanerVM.Arithmetization.Caps.blake2s_height : 2 ^ minLogRowsBlake2s ≤ (blake2sRows w).length
  LeanerVM.Arithmetization.Caps.well_shaped : WellShapedData w.data
constructor:
  LeanerVM.Arithmetization.Caps.mk {prog : Program} {w : EnsembleWitness (leanIsaEnsemble prog)}
    (minLogMem_le : minLogMem ≤ (imageOf w.data).fst) (le_maxLogMem : (imageOf w.data).fst ≤ maxLogMem)
    (heights : ∀ t ∈ w.tables, ∃ τ ≤ maxLogRows, t.table.length = 2 ^ τ)
    (blake2s_height : 2 ^ minLogRowsBlake2s ≤ (blake2sRows w).length) (well_shaped : WellShapedData w.data) : Caps w
structure LeanerVM.Semantics.PublicInput : Type
number of parameters: 0
fields:
  LeanerVM.Semantics.PublicInput.lanes : Fin 4 → K
constructor:
  LeanerVM.Semantics.PublicInput.mk (lanes : Fin 4 → K) : PublicInput
def LeanerVM.Semantics.PublicInput.word0 : PublicInput → E :=
fun p => E.ofLimbs (p.lanes 0) (p.lanes 1) 0
def LeanerVM.Semantics.PublicInput.word1 : PublicInput → E :=
fun p => E.ofLimbs (p.lanes 2) (p.lanes 3) 0
structure LeanerVM.Arithmetization.PublicIO (F : Type) : Type
number of parameters: 1
fields:
  LeanerVM.Arithmetization.PublicIO.lanes : Vector F 4
constructor:
  LeanerVM.Arithmetization.PublicIO.mk {F : Type} (lanes : Vector F 4) : PublicIO F
def LeanerVM.Semantics.ValidExecution : (prog : Program) → PublicInput → Trace prog → Prop :=
fun prog input t => HasPublicBoundary input t ∧ run prog t.image t.steps Regs.initial = some (Regs.final prog)
def LeanerVM.Semantics.HasPublicBoundary : {prog : Program} → PublicInput → Trace prog → Prop :=
fun {prog} input t =>
  minLogMem ≤ t.κ ∧
    t.κ ≤ maxLogMem ∧ t.image.read (gpow 0) = some input.word0 ∧ t.image.read (gpow 1) = some input.word1
structure LeanerVM.Semantics.Program : Type
number of parameters: 0
fields:
  LeanerVM.Semantics.Program.logSize : ℕ
  LeanerVM.Semantics.Program.logSize_le : self.logSize ≤ maxLogBytecode
  LeanerVM.Semantics.Program.code : Fin (2 ^ self.logSize) → Instr
constructor:
  LeanerVM.Semantics.Program.mk (logSize : ℕ) (logSize_le : logSize ≤ maxLogBytecode)
    (code : Fin (2 ^ logSize) → Instr) : Program
def LeanerVM.Protocol.M3Holds : (I : M3Instance) → I.Stmt → Column I.μ → Prop :=
fun I input q => I.ConstraintsVanish q ∧ I.Balanced q ∧ I.CountsNonzero q ∧ I.PublicLinesHold input q ∧ I.aux q
structure LeanerVM.Protocol.M3Instance : Type 1
number of parameters: 0
parents:
  LeanerVM.Protocol.M3Instance.toShape : Shape
fields:
  LeanerVM.Protocol.Shape.ntab : ℕ
  LeanerVM.Protocol.Shape.τ : Fin self.ntab → ℕ
  LeanerVM.Protocol.Shape.width : Fin self.ntab → ℕ
  LeanerVM.Protocol.M3Instance.Stmt : Type
  LeanerVM.Protocol.M3Instance.constraints : (j : Fin self.ntab) → List (CPoly.CMvPolynomial (self.width j) K)
  LeanerVM.Protocol.M3Instance.flushes : (j : Fin self.ntab) →
      List (Side × Vector (CPoly.CMvPolynomial (self.width j) K) 16)
  LeanerVM.Protocol.M3Instance.d : ℕ
  LeanerVM.Protocol.M3Instance.constraints_degree : ∀ (j : Fin self.ntab),
      ∀ C ∈ self.constraints j, C.totalDegree ≤ self.d
  LeanerVM.Protocol.M3Instance.flushes_degree : ∀ (j : Fin self.ntab),
      ∀ f ∈ self.flushes j, ∀ (k : Fin 16), (f.2.get k).totalDegree ≤ self.d
  LeanerVM.Protocol.M3Instance.counts : (j : Fin self.ntab) → List (Fin (self.width j))
  LeanerVM.Protocol.M3Instance.boundary : List (BoundaryBlock self.toShape)
  LeanerVM.Protocol.M3Instance.μ : ℕ
  LeanerVM.Protocol.M3Instance.layout : Layout self.μ self.ColumnId fun c => self.τ c.fst
  LeanerVM.Protocol.M3Instance.publicLines : self.Stmt → List (PublicLine self.toShape)
  LeanerVM.Protocol.M3Instance.aux : Column self.μ → Prop
  LeanerVM.Protocol.M3Instance.decAux : DecidablePred self.aux
constructor:
  LeanerVM.Protocol.M3Instance.mk (toShape : Shape) (Stmt : Type)
    (constraints : (j : Fin toShape.ntab) → List (CPoly.CMvPolynomial (toShape.width j) K))
    (flushes : (j : Fin toShape.ntab) → List (Side × Vector (CPoly.CMvPolynomial (toShape.width j) K) 16)) (d : ℕ)
    (constraints_degree : ∀ (j : Fin toShape.ntab), ∀ C ∈ constraints j, C.totalDegree ≤ d)
    (flushes_degree : ∀ (j : Fin toShape.ntab), ∀ f ∈ flushes j, ∀ (k : Fin 16), (f.2.get k).totalDegree ≤ d)
    (counts : (j : Fin toShape.ntab) → List (Fin (toShape.width j))) (boundary : List (BoundaryBlock toShape)) (μ : ℕ)
    (layout : Layout μ toShape.ColumnId fun c => toShape.τ c.fst) (publicLines : Stmt → List (PublicLine toShape))
    (aux : Column μ → Prop) (decAux : DecidablePred aux) : M3Instance
field notation resolution order:
  LeanerVM.Protocol.M3Instance, LeanerVM.Protocol.Shape
structure LeanerVM.Protocol.PublicLine (S : Shape) : Type
number of parameters: 1
fields:
  LeanerVM.Protocol.PublicLine.col : S.ColumnId
  LeanerVM.Protocol.PublicLine.cell0 : K
  LeanerVM.Protocol.PublicLine.cell1 : K
  LeanerVM.Protocol.PublicLine.sent : Bool
  LeanerVM.Protocol.PublicLine.pos : 0 < S.τ self.col.fst
constructor:
  LeanerVM.Protocol.PublicLine.mk {S : Shape} (col : S.ColumnId) (cell0 cell1 : K) (sent : Bool)
    (pos : 0 < S.τ col.fst) : PublicLine S
def Air.Flat.EnsembleWitness.Constraints : {F : Type} →
  [inst : FiniteField F] →
    {PublicIO : TypeMap} →
      [inst_1 : ProvableType PublicIO] → {ens : Ensemble F PublicIO} → EnsembleWitness ens → Prop :=
fun {F} [FiniteField F] {PublicIO} [ProvableType PublicIO] {ens} witness =>
  ∀ table ∈ witness.allTables, table.Constraints
def Air.Flat.Table.Constraints : {F : Type} → [inst : FiniteField F] → Air.Flat.Table F → Prop :=
fun {F} [FiniteField F] table =>
  ∀ row ∈ table.table, Operations.ConstraintsHold (table.environment row) table.component.operations
def Operations.ConstraintsHold : {F : Type} → [inst : FiniteField F] → Environment F → Operations F → Prop :=
fun {F} [FiniteField F] env ops =>
  (∀ e ∈ ops.constraints, (fun x => Expression.eval env x) e = 0) ∧ ∀ l ∈ ops.lookups, l.Contains env
def Air.Flat.EnsembleWitness.interactions : {F : Type} →
  [inst : FiniteField F] →
    {PublicIO : TypeMap} →
      [inst_1 : ProvableType PublicIO] → {ens : Ensemble F PublicIO} → EnsembleWitness ens → List (Interaction F) :=
fun {F} [FiniteField F] {PublicIO} [ProvableType PublicIO] {ens} witness =>
  List.flatMap (fun table => table.interactions) witness.allTables
def ProverData : Type → Type :=
fun F => String → (n : ℕ) → Array (Vector F n)
'LeanerVM.Arithmetization.assumptions_of_blake2sRowsValid' depends on axioms: [propext, Classical.choice, Quot.sound]
'LeanerVM.Arithmetization.memRowOf_bindings' depends on axioms: [propext, Classical.choice, Quot.sound]
'LeanerVM.Arithmetization.bytecodeRowOf_bindings' depends on axioms: [propext, Classical.choice, Quot.sound]
'LeanerVM.Arithmetization.bytecodeRowOf_decodes' depends on axioms: [propext, Classical.choice, Quot.sound]
'LeanerVM.Arithmetization.verifier_pull_eval' depends on axioms: [propext, Classical.choice, Quot.sound]
'LeanerVM.Arithmetization.verifier_push_eval' depends on axioms: [propext, Classical.choice, Quot.sound]
'LeanerVM.Arithmetization.rowAt_toElements' depends on axioms: [propext, Classical.choice, Quot.sound]
'LeanerVM.Arithmetization.decode_eq_some_iff' depends on axioms: [propext, Classical.choice, Quot.sound]
'LeanerVM.Arithmetization.assignmentRepresents_image' depends on axioms: [propext, Classical.choice, Quot.sound]
'LeanerVM.Arithmetization.xorTable' depends on axioms: [propext, Classical.choice, Quot.sound]
'LeanerVM.Arithmetization.jumpTable' depends on axioms: [propext, Classical.choice, Quot.sound]
'LeanerVM.Arithmetization.blake2sTable' depends on axioms: [propext, Classical.choice, Quot.sound]
'LeanerVM.Arithmetization.memTable' depends on axioms: [propext, Classical.choice, Quot.sound]
'LeanerVM.Arithmetization.bytecodeTable' depends on axioms: [propext, Classical.choice, Quot.sound]
'LeanerVM.Arithmetization.leanIsaVerifier' depends on axioms: [propext, Classical.choice, Quot.sound]
'LeanerVM.Protocol.piop_rbrKnowledgeSoundness' depends on axioms: [propext, Classical.choice, Quot.sound]
'LeanerVM.Protocol.piop_perfectCompleteness' depends on axioms: [propext, Classical.choice, Quot.sound]
'LeanerVM.Protocol.Refinement.map_option_valid' depends on axioms: [propext, Quot.sound]
'LeanerVM.Protocol.bytecodeColumn_eval' depends on axioms: [propext, Classical.choice, Quot.sound]
'LeanerVM.Protocol.idxColumn_eval' depends on axioms: [propext, Classical.choice, Quot.sound]
exit=0
\end{srccode}

\subsection{\texorpdfstring{\protect\code{KnownColumn}}{KnownColumn}}\label{sec:gen-boundary-probe-knowncolumn}

\paragraph{What it establishes.} The public column of a boundary block is part of the
relation, and an instance that treats it as committed while the adaptor rebuilds it falls into
the trap of \cref{sec:gen-boundary-trap}. A model on the spine's toy instance: \bndc{toyOf p} has
one boundary block whose coordinate 0 is \bndc{Coord.known p}, the column \bndc{p} standing for
the program; \bndc{toyCommitted} reads the same coordinate off column 1 of the table, as if the
program were committed. The honest stack satisfies \bndc{M3Holds (toyOf prog)} and fails
\bndc{Balanced} for another program, so the relation is about the program in the instance. The
stack \bndc{pushesTwo} satisfies \bndc{M3Holds toyCommitted} while \bndc{(toyOf prog).Balanced}
fails on it: the witness an adaptor would rebuild from it, the block built from the program, does
not balance. The numeral \bndc{2 : K} is nonzero at the old CompPoly pin, as the spine's own test
\bndc{badConstraint} uses it; at the new pin it is \bndc{0}, hence the variant
\bndc{KnownColumnNew}.

\paragraph{Revision.} Run on 2026-09-29 against leanerVM \code{main} at \code{b435631} with the old pins (ArkLib \code{dca90385}, CompPoly \code{3468b38c}, Clean \code{93c9d1ef}). Exit 0.

\paragraph{Source.}
\begin{leancode}[numbers=left,numbersep=3pt]
/-
  Probe (boundary-adaptor): the public column of a boundary block is part of the relation, and
  what goes wrong when an instance treats it as committed while the adaptor rebuilds it.
  A model on the spine's toy instance; scratch work for the blueprint review.
  Run from the repository root:
    flock .claude/reports/blueprint-review/logs/lean.lock \
      lake env lean .claude/reports/blueprint-review/probes/boundary-adaptor/KnownColumn.lean
-/
import LeanerVM.Protocol.Spine.Toy

open LeanerVM.Parameters LeanerVM.Protocol LeanerVM.Protocol.Toy

namespace Probe

/-- The toy instance with the boundary block's public column a parameter: the model of
`leanIsaInstance prog s`, the column `p` standing for the program. -/
abbrev toyOf (p : Column 1) : M3Instance :=
  { toy with
    boundary := [{ κ := 1, side := .pull,
                   coords := Vector.ofFn fun k ↦ if k.val = 0 then .known p else .const 0 }] }

/-- The same instance built wrongly: the block's column is read from the stack (column 1 of the
table), as if the program were committed. -/
abbrev toyCommitted : M3Instance :=
  { toy with
    boundary := [{ κ := 1, side := .pull,
                   coords := Vector.ofFn fun k ↦
                     if k.val = 0 then .committed ⟨0, 1⟩ rfl else .const 0 }] }

/-- The public program of the model. -/
def prog : Column 1 := ⟨#v[1, 1]⟩

/-- Another program. -/
def prog' : Column 1 := ⟨#v[1, 2]⟩

-- 1. The relation is about the program in the instance: the honest stack of `prog` satisfies the
-- relation of `prog`, and not the relation of `prog'`.
#guard M3Holds (toyOf prog) (1 : K) honest
#guard ¬ (toyOf prog').Balanced honest
#guard ¬ M3Holds (toyOf prog') (1 : K) honest

/-- A stack that pushes `[2, 2]` and whose column 1 is `[2, 2]`. -/
def pushesTwo : Column 3 := ⟨#v[2, 2, 2, 2, 1, 0, 0, 0]⟩

-- 2. The trap. Against the instance that reads the block's column off the stack, the relation
-- holds of a stack that has nothing to do with the public program ...
#guard M3Holds toyCommitted (1 : K) pushesTwo
-- ... while the witness an adaptor would rebuild from it, the table's columns read from the
-- stack and the block's column built from the public program, does not balance.
#guard ¬ (toyOf prog).Balanced pushesTwo
-- The pulled tuples of the two instances on the same stack: the first from the stack, the second
-- from the program.
#guard (toyCommitted.tuples pushesTwo .pull).map (fun t ↦ t.get 0) = [2, 2]
#guard ((toyOf prog).tuples pushesTwo .pull).map (fun t ↦ t.get 0) = [1, 1]

end Probe
\end{leancode}
\src{.claude/reports/blueprint-review/probes/boundary-adaptor/KnownColumn.lean}

\paragraph{Command} (from the repository root).
\begin{shellcode}
flock .claude/reports/blueprint-review/logs/lean.lock \
  lake env lean .claude/reports/blueprint-review/probes/boundary-adaptor/KnownColumn.lean
\end{shellcode}

\paragraph{Recorded output} (\file{KnownColumn.out}).
\begin{srccode}[numbers=left,numbersep=3pt]
exit=0
\end{srccode}

\subsection{\texorpdfstring{\protect\code{KnownColumnNew}}{KnownColumnNew}}\label{sec:gen-boundary-probe-knowncolumnnew}

\paragraph{What it establishes.} The same as \bndc{KnownColumn}
(\cref{sec:gen-boundary-probe-knowncolumn}) at the new pins, with \bndc{K.ofBits 2} in place of the
numeral \bndc{2}; every guard passes.

\paragraph{Revision.} Run on 2026-09-30, after the rebuild, against the checkout \code{docs/protocol-blueprint-review} at \code{b022e4f}, whose Lean sources are \code{main} at \code{144c5aa}, with the new pins (CompPoly \code{572f9973}, Clean \code{42fe4b26}, ArkLib \code{7653a901}). Exit 0.

\paragraph{Source.}
\begin{leancode}[numbers=left,numbersep=3pt]
/-
  Probe (boundary-adaptor), NEW-PIN VARIANT (CompPoly 572f9973, K.ofBits in place of numerals): the public column of a boundary block is part of the relation, and
  what goes wrong when an instance treats it as committed while the adaptor rebuilds it.
  A model on the spine's toy instance; scratch work for the blueprint review.
  Run from the repository root:
    flock .claude/reports/blueprint-review/logs/lean.lock \
      lake env lean .claude/reports/blueprint-review/probes/boundary-adaptor/KnownColumnNew.lean
-/
import LeanerVM.Protocol.Spine.Toy

open LeanerVM.Parameters LeanerVM.Protocol LeanerVM.Protocol.Toy

namespace Probe

/-- The toy instance with the boundary block's public column a parameter: the model of
`leanIsaInstance prog s`, the column `p` standing for the program. -/
abbrev toyOf (p : Column 1) : M3Instance :=
  { toy with
    boundary := [{ κ := 1, side := .pull,
                   coords := Vector.ofFn fun k ↦ if k.val = 0 then .known p else .const 0 }] }

/-- The same instance built wrongly: the block's column is read from the stack (column 1 of the
table), as if the program were committed. -/
abbrev toyCommitted : M3Instance :=
  { toy with
    boundary := [{ κ := 1, side := .pull,
                   coords := Vector.ofFn fun k ↦
                     if k.val = 0 then .committed ⟨0, 1⟩ rfl else .const 0 }] }

/-- The public program of the model. -/
def prog : Column 1 := ⟨#v[1, 1]⟩

/-- Another program. -/
def prog' : Column 1 := ⟨#v[1, K.ofBits 2]⟩

-- 1. The relation is about the program in the instance: the honest stack of `prog` satisfies the
-- relation of `prog`, and not the relation of `prog'`.
#guard M3Holds (toyOf prog) (1 : K) honest
#guard ¬ (toyOf prog').Balanced honest
#guard ¬ M3Holds (toyOf prog') (1 : K) honest

/-- A stack that pushes `[2, 2]` and whose column 1 is `[2, 2]`. -/
def pushesTwo : Column 3 := ⟨#v[K.ofBits 2, K.ofBits 2, K.ofBits 2, K.ofBits 2, 1, 0, 0, 0]⟩

-- 2. The trap. Against the instance that reads the block's column off the stack, the relation
-- holds of a stack that has nothing to do with the public program ...
#guard M3Holds toyCommitted (1 : K) pushesTwo
-- ... while the witness an adaptor would rebuild from it, the table's columns read from the
-- stack and the block's column built from the public program, does not balance.
#guard ¬ (toyOf prog).Balanced pushesTwo
-- The pulled tuples of the two instances on the same stack: the first from the stack, the second
-- from the program.
#guard (toyCommitted.tuples pushesTwo .pull).map (fun t ↦ t.get 0) = [K.ofBits 2, K.ofBits 2]
#guard ((toyOf prog).tuples pushesTwo .pull).map (fun t ↦ t.get 0) = [1, 1]

end Probe
\end{leancode}
\src{.claude/reports/blueprint-review/probes/boundary-adaptor/KnownColumnNew.lean}

\paragraph{Command} (from the repository root).
\begin{shellcode}
flock .claude/reports/blueprint-review/logs/lean.lock \
  lake env lean .claude/reports/blueprint-review/probes/boundary-adaptor/KnownColumnNew.lean
\end{shellcode}

\paragraph{Recorded output} (\file{KnownColumnNew.out}).
\begin{srccode}[numbers=left,numbersep=3pt]
exit=0
\end{srccode}

\subsection{\texorpdfstring{\protect\code{Transport}}{Transport}}\label{sec:gen-boundary-probe-transport}

\paragraph{What it establishes.} Whether the chain \enquote{knowledge of \bndc{M3Holds}
$→$ \bndc{SatisfiedBy} $→$ \bndc{ValidExecution}} typechecks with the adaptor's target witness in
\bndc{Type 1}. The instance, the witness map and the two theorems not yet built are hypotheses of
each \bndc{example}; nothing is an axiom. Example 1: the adaptor is a \bndc{Refinement}, which
accepts a target witness in \bndc{Type 1}. \bndc{bad_of_bad}: the transport of the bad event of
ArkLib's game (\enquote{no witness in the slot is valid}), which follows from \bndc{map_valid} in
one line, where \bndc{Refinement.map_option_valid} has the other polarity. Examples 3 and 4: the
pointwise composition; when \bndc{prog} has no valid execution on \bndc{input}, whatever the
extractor returns is outside \bndc{M3Rel}, so the event of ArkLib's game is \enquote{the verifier
accepts}, and the conclusion of the chain is a soundness statement for the language
\bndc{{(prog, input) | ∃ t, ValidExecution prog input t}}; with the hypothesis
\bndc{SentinelSafe prog} of constraint soundness, the hypothesis reaches the conclusion. Used in
\cref{sec:gen-boundary-st-composition}.

\paragraph{Revision.} Run on 2026-09-29 against leanerVM \code{main} at \code{b435631} with the old pins (ArkLib \code{dca90385}, CompPoly \code{3468b38c}, Clean \code{93c9d1ef}). Exit 0. Not re-run at the new pins.

\paragraph{Source.}
\begin{leancode}[numbers=left,numbersep=3pt]
/-
  Probe (boundary-adaptor): does the chain
    knowledge of `M3Holds`  →  `SatisfiedBy`  →  `ValidExecution`
  typecheck, with the adaptor's target witness in `Type 1`?
  The instance, the witness map and the two theorems not yet built are hypotheses of each
  `example`; nothing here is an axiom. Scratch work for the blueprint review.
  Run from the repository root:
    flock .claude/reports/blueprint-review/logs/lean.lock \
      lake env lean .claude/reports/blueprint-review/probes/boundary-adaptor/Transport.lean
-/
import LeanerVM

open LeanerVM LeanerVM.Parameters LeanerVM.Semantics LeanerVM.Arithmetization
open LeanerVM.Protocol
open Air.Flat (Ensemble EnsembleWitness Component)

namespace Probe

/-- A stand-in for `leanIsaInstance prog s`: any instance whose statement type is
`PublicInput`. (The toy's tables, no public line.) -/
abbrev stub : M3Instance :=
  { Toy.toy with Stmt := PublicInput, publicLines := fun _ ↦ [] }

/-- The relation of leanISA in ArkLib's shape, over the statement type of `M3Rel`. Its witness
type is in `Type 1`. -/
def SatRel (prog : Program) :
    Set ((PublicInput × (∀ i, NoOracle i)) × EnsembleWitness (leanIsaEnsemble prog)) :=
  {p | SatisfiedBy prog p.1.1 p.2}

/-- 1. The adaptor is a `Refinement`: the structure accepts a target witness in `Type 1`. -/
example (prog : Program) (witnessOf : Column stub.μ → EnsembleWitness (leanIsaEnsemble prog))
    (satisfiedBy_witnessOf :
      ∀ input q, M3Holds stub input q → SatisfiedBy prog input (witnessOf q)) :
    Refinement (M3Rel stub) (SatRel prog) where
  map := fun _ q ↦ witnessOf q
  map_valid := fun x q hx ↦ satisfiedBy_witnessOf x.1 q hx

/-- 2. The polarity ArkLib's game needs. Its bad event is "no witness in the slot is valid"
(an empty slot is bad); `Refinement.map_option_valid` is stated for "every witness in the slot
is valid" (an empty slot is good). The transport of the bad event is this lemma, which follows
from `map_valid` directly. -/
theorem bad_of_bad {Stmt : Type} {W₁ : Type} {W₂ : Type 1} {R : Set (Stmt × W₁)}
    {S : Set (Stmt × W₂)} (f : Refinement R S) (x : Stmt) (w? : Option W₁)
    (h : ∀ w' ∈ w?.map (f.map x), (x, w') ∉ S) : ∀ w ∈ w?, (x, w) ∉ R := by
  intro w hw hR
  exact h (f.map x w) (Option.mem_map_of_mem _ hw) (f.map_valid x w hR)

/-- 3. The composition of the chain, pointwise, in the form the probability bound consumes:
when `prog` has no valid execution on `input`, whatever the extractor returns is outside
`M3Rel`, so the event of ArkLib's knowledge-soundness game is the event "the verifier accepts".
The conclusion of the chain is therefore a soundness statement for the language
`{(prog, input) | ∃ t, ValidExecution prog input t}`. -/
example (prog : Program) (input : PublicInput)
    (witnessOf : Column stub.μ → EnsembleWitness (leanIsaEnsemble prog))
    (satisfiedBy_witnessOf :
      ∀ q, M3Holds stub input q → SatisfiedBy prog input (witnessOf q))
    (constraintSoundness :
      ∀ w, SatisfiedBy prog input w → ∃ t, ValidExecution prog input t)
    (hno : ¬ ∃ t, ValidExecution prog input t) (o : ∀ i, NoOracle i)
    (q? : Option (Column stub.μ)) :
    ∀ q ∈ q?, ((input, o), q) ∉ M3Rel stub := by
  intro q _ hq
  exact hno (constraintSoundness _ (satisfiedBy_witnessOf q hq))

/-- 4. The same with the hypothesis of `constraintSoundness` the leanISA roadmap states
(`WellFormedBytecode prog`, of which soundness uses the sentinel field, built as
`SentinelSafe prog`): the hypothesis reaches the conclusion of the chain. -/
example (prog : Program) (input : PublicInput)
    (witnessOf : Column stub.μ → EnsembleWitness (leanIsaEnsemble prog))
    (satisfiedBy_witnessOf :
      ∀ q, M3Holds stub input q → SatisfiedBy prog input (witnessOf q))
    (constraintSoundness :
      SentinelSafe prog → ∀ w, SatisfiedBy prog input w → ∃ t, ValidExecution prog input t)
    (hwf : SentinelSafe prog)
    (q : Column stub.μ) (h : M3Holds stub input q) : ∃ t, ValidExecution prog input t :=
  constraintSoundness hwf _ (satisfiedBy_witnessOf q h)

end Probe
\end{leancode}
\src{.claude/reports/blueprint-review/probes/boundary-adaptor/Transport.lean}

\paragraph{Command} (from the repository root).
\begin{shellcode}
flock .claude/reports/blueprint-review/logs/lean.lock \
  lake env lean .claude/reports/blueprint-review/probes/boundary-adaptor/Transport.lean
\end{shellcode}

\paragraph{Recorded output} (\file{Transport.out}).
\begin{srccode}[numbers=left,numbersep=3pt]
exit=0
\end{srccode}

\subsection{\texorpdfstring{\protect\code{UniverseFail}}{UniverseFail}}\label{sec:gen-boundary-probe-universefail}

\paragraph{What it establishes.} Expected to fail. ArkLib's straight-line extractor
cannot be stated with the witness of \bndc{SatisfiedBy}, which lives in \bndc{Type 1}: both an
extractor type whose witness is an \bndc{EnsembleWitness} and the post-composition
\bndc{Extractor.Straightline.map witnessOf E} are rejected with a universe mismatch. So
\bndc{Extractor.Straightline.map} cannot serve the adaptor
(\cref{sec:gen-boundary-st-composition}).

\paragraph{Revision.} Run on 2026-09-29 against leanerVM \code{main} at \code{b435631} with the old pins (ArkLib \code{dca90385}, CompPoly \code{3468b38c}, Clean \code{93c9d1ef}). Exit 1, with the two recorded errors, as expected. Not re-run at the new pins.

\paragraph{Source.}
\begin{leancode}[numbers=left,numbersep=3pt]
/-
  Probe (boundary-adaptor), EXPECTED TO FAIL: ArkLib's knowledge-soundness game and its
  straight-line extractor cannot be stated with the witness of `SatisfiedBy`, which lives in
  `Type 1`. Scratch work for the blueprint review.
  Run from the repository root:
    flock .claude/reports/blueprint-review/logs/lean.lock \
      lake env lean .claude/reports/blueprint-review/probes/boundary-adaptor/UniverseFail.lean
-/
import LeanerVM

open LeanerVM LeanerVM.Parameters LeanerVM.Semantics LeanerVM.Arithmetization
open LeanerVM.Protocol OracleComp OracleSpec ProtocolSpec
open Air.Flat (Ensemble EnsembleWitness Component)

namespace Probe

/-- Expected failure 1: a straight-line extractor whose output is an `EnsembleWitness`. -/
example (prog : Program) {n : ℕ} (pSpec : ProtocolSpec n) : Type :=
  Extractor.Straightline []ₒ PublicInput (EnsembleWitness (leanIsaEnsemble prog)) Unit pSpec

/-- Expected failure 2: post-composing an extractor of stacks with a map to `EnsembleWitness`. -/
example (prog : Program) {n : ℕ} (pSpec : ProtocolSpec n) (μ : ℕ)
    (witnessOf : PublicInput → Column μ → EnsembleWitness (leanIsaEnsemble prog))
    (E : Extractor.Straightline []ₒ PublicInput (Column μ) Unit pSpec) :=
  Extractor.Straightline.map witnessOf E

end Probe
\end{leancode}
\src{.claude/reports/blueprint-review/probes/boundary-adaptor/UniverseFail.lean}

\paragraph{Command} (from the repository root).
\begin{shellcode}
flock .claude/reports/blueprint-review/logs/lean.lock \
  lake env lean .claude/reports/blueprint-review/probes/boundary-adaptor/UniverseFail.lean
\end{shellcode}

\paragraph{Recorded output} (\file{UniverseFail.out}).
\begin{srccode}[numbers=left,numbersep=3pt]
.claude/reports/blueprint-review/probes/boundary-adaptor/UniverseFail.lean:19:41: error: Application type mismatch: The argument
  EnsembleWitness (leanIsaEnsemble prog)
has type
  Type 1
of sort `Type 2` but is expected to have type
  Type
of sort `Type 1` in the application
  Extractor.Straightline []ₒ PublicInput (EnsembleWitness (leanIsaEnsemble prog))
.claude/reports/blueprint-review/probes/boundary-adaptor/UniverseFail.lean:25:29: error: Application type mismatch: The argument
  witnessOf
has type
  PublicInput → Column μ → EnsembleWitness (leanIsaEnsemble prog)
of sort `Type 1` but is expected to have type
  PublicInput → Column μ → ?m.27 prog pSpec μ witnessOf E
of sort `Type` in the application
  Extractor.Straightline.map witnessOf
exit=1
\end{srccode}

\subsection{\texorpdfstring{\protect\code{DefInstance}}{DefInstance}}\label{sec:gen-boundary-probe-definstance}

\paragraph{What it establishes.} Expected to fail in part. With a stand-in for
\bndc{leanIsaInstance prog s} written as a \bndc{def} with parameters, the statement of
\bndc{satisfiedBy_witnessOf}'s hypothesis elaborates (the \bndc{example} at lines 33--34: the
elaborator unfolds the definition to see the statement type), but instance search does not see
through it: the two \bndc{#guard}s on \bndc{inst} fail with \enquote{Type mismatch … expected to
have type Bool}, no \bndc{Decidable} instance being found, while the same guards on the
\bndc{abbrev} \bndc{instA} pass. Used in \cref{sec:gen-boundary-st-witnessof} and
\cref{f:boundary-reducible}.

\paragraph{Revision.} Run on 2026-09-29 against leanerVM \code{main} at \code{b435631} with the old pins (ArkLib \code{dca90385}, CompPoly \code{3468b38c}, Clean \code{93c9d1ef}). Exit 1, with the recorded errors on the two guards of the \code{def}, as expected. Not re-run at the new pins.

\paragraph{Source.}
\begin{leancode}[numbers=left,numbersep=3pt]
/-
  Probe (boundary-adaptor): does the adaptor's statement typecheck when the instance is a `def`
  (not an `abbrev`), and is its relation still decided by evaluation?
  The two `#guard`s on the `def` instance are EXPECTED TO FAIL (no `Decidable` instance is
  found); the same guards on the `abbrev` instance pass.
  Scratch work for the blueprint review.
  Run from the repository root:
    flock .claude/reports/blueprint-review/logs/lean.lock \
      lake env lean .claude/reports/blueprint-review/probes/boundary-adaptor/DefInstance.lean
-/
import LeanerVM.Protocol.Spine.Toy
import LeanerVM.Semantics.Instruction

open LeanerVM.Parameters LeanerVM.Semantics LeanerVM.Protocol LeanerVM.Protocol.Toy

namespace Probe

/-- A stand-in for `leanIsaInstance prog s`, as a `def` with parameters: the statement type is
`PublicInput`, and the one public line reads lane 0 of the input. -/
def inst (_prog : Program) (_s : ℕ) : M3Instance :=
  { toy with
    Stmt := PublicInput
    publicLines := fun input ↦ [⟨⟨0, 2⟩, input.lanes 0, 0, true, by decide⟩] }

/-- The same, as an `abbrev`. -/
abbrev instA (_prog : Program) (_s : ℕ) : M3Instance :=
  { toy with
    Stmt := PublicInput
    publicLines := fun input ↦ [⟨⟨0, 2⟩, input.lanes 0, 0, true, by decide⟩] }

/-- The statement of `satisfiedBy_witnessOf`'s hypothesis elaborates with `input : PublicInput`
on the `def`: the elaborator unfolds it to see the statement type. -/
example (prog : Program) (s : ℕ) (input : PublicInput) (q : Column (inst prog s).μ) : Prop :=
  M3Holds (inst prog s) input q

/-- A program and an input for the evaluation below. -/
def prog0 : Program := ⟨0, by decide, fun _ ↦ .xor 1 1 1⟩
def input1 : PublicInput := ⟨fun _ ↦ 1⟩
def input0 : PublicInput := ⟨fun _ ↦ 0⟩

-- On the `abbrev`, the relation is decided by evaluation, as the toy's is.
#guard M3Holds (instA prog0 0) input1 honest
#guard ¬ M3Holds (instA prog0 0) input0 honest

-- On the `def`, instance search does not see through the definition (EXPECTED TO FAIL).
#guard M3Holds (inst prog0 0) input1 honest
#guard ¬ M3Holds (inst prog0 0) input0 honest

end Probe
\end{leancode}
\src{.claude/reports/blueprint-review/probes/boundary-adaptor/DefInstance.lean}

\paragraph{Command} (from the repository root).
\begin{shellcode}
flock .claude/reports/blueprint-review/logs/lean.lock \
  lake env lean .claude/reports/blueprint-review/probes/boundary-adaptor/DefInstance.lean
\end{shellcode}

\paragraph{Recorded output} (\file{DefInstance.out}).
\begin{srccode}[numbers=left,numbersep=3pt]
.claude/reports/blueprint-review/probes/boundary-adaptor/DefInstance.lean:46:7: error: Type mismatch
  M3Holds (inst prog0 0) input1 honest
has type
  Prop
but is expected to have type
  Bool
.claude/reports/blueprint-review/probes/boundary-adaptor/DefInstance.lean:46:0: error: cannot evaluate code because 'sorryAx' uses 'sorry' and/or contains errors
.claude/reports/blueprint-review/probes/boundary-adaptor/DefInstance.lean:47:7: error: Type mismatch
  ¬M3Holds (inst prog0 0) input0 honest
has type
  Prop
but is expected to have type
  Bool
.claude/reports/blueprint-review/probes/boundary-adaptor/DefInstance.lean:47:0: error: cannot evaluate code because 'sorryAx' uses 'sorry' and/or contains errors
exit=1
\end{srccode}

\subsection{\texorpdfstring{\protect\code{PolyBridge}}{PolyBridge}}\label{sec:gen-boundary-probe-polybridge}

\paragraph{What it establishes.} From a Clean expression to the polynomial an
\bndc{M3Instance} holds. It prints CompPoly's two conversions at the pin, then defines
\bndc{exprToCMv}, a direct translation of a Clean expression over \bndc{K} to a computable
\bndc{CMvPolynomial} in twelve lines, a variable past the width reading \bndc{0} as
\bndc{Environment.fromArray} reads a missing cell; and checks by \bndc{#guard} that it agrees with
Clean's \bndc{Expression.eval} on a row for the first \bndc{JUMP} residual \bndc{b + v_cond · w},
that the polynomial's \bndc{totalDegree} is 2, and that a variable past the width reads \bndc{0} on
both sides (\cref{sec:gen-boundary-st-poly}). That \bndc{toCMvPolynomial} is \bndc{noncomputable}
rests on its source (\bndc{Multivariate/MvPolyEquiv/Core.lean:41}). At the new CompPoly pin the
row's numerals \bndc{3, 5, 7} read as \bndc{1, 1, 1}, hence the variant \bndc{PolyBridgeNew}.

\paragraph{Revision.} Run on 2026-09-29 against leanerVM \code{main} at \code{b435631} with the old pins (ArkLib \code{dca90385}, CompPoly \code{3468b38c}, Clean \code{93c9d1ef}). Exit 0.

\paragraph{Source.}
\begin{leancode}[numbers=left,numbersep=3pt]
/-
  Probe (boundary-adaptor): from a Clean expression to the polynomial an `M3Instance` holds.
  The blueprint's Layer 2 produces Mathlib's `MvPolynomial`; the instance holds CompPoly's
  `CMvPolynomial`. CompPoly's conversion from the first to the second is noncomputable at the
  pin; a direct, computable translation exists. Scratch work for the blueprint review.
  Run from the repository root:
    flock .claude/reports/blueprint-review/logs/lean.lock \
      lake env lean .claude/reports/blueprint-review/probes/boundary-adaptor/PolyBridge.lean
-/
import LeanerVM
import CompPoly.Multivariate.MvPolyEquiv.Eval
import CompPoly.Multivariate.Operations

open LeanerVM.Parameters CompPoly CPoly

namespace Probe

-- The two directions of CompPoly's bridge at the pin.
#print CPoly.toCMvPolynomial
#print CPoly.fromCMvPolynomial

/-- A direct translation of a Clean expression over `K` to a computable polynomial in the first
`n` row variables; a variable past the width reads `0`, as Clean's `Environment.fromArray`
reads a missing cell. -/
def exprToCMv (n : ℕ) : Expression K → CMvPolynomial n K
  | .var v => if h : v.index < n then CMvPolynomial.X ⟨v.index, h⟩ else 0
  | .const c => CMvPolynomial.C c
  | .add a b => exprToCMv n a + exprToCMv n b
  | .mul a b => exprToCMv n a * exprToCMv n b

/-- The first `JUMP` residual, `b + v_cond · w`, on a row `(v_cond, w, b)`. -/
def residual : Expression K := .add (.var ⟨2⟩) (.mul (.var ⟨0⟩) (.var ⟨1⟩))

/-- A row. -/
def row : Array K := #[3, 5, 7]

/-- The row as the assignment of the three variables. -/
def rowFn : Fin 3 → K := fun i ↦ row[i.val]?.getD 0

/-- Clean's value of the residual on the row. -/
def cleanValue : K := residual.eval (Environment.fromArray row fun _ _ ↦ #[])

/-- The polynomial's value on the row. -/
def polyValue : K := (exprToCMv 3 residual).eval rowFn

-- The translation runs, agrees with Clean's evaluation on the row, and has the degree the
-- expression shows.
#guard polyValue = cleanValue
#guard (exprToCMv 3 residual).totalDegree = 2

-- A variable past the width is read as zero on both sides.
def wide : Expression K := .add (.var ⟨5⟩) (.var ⟨0⟩)
def wideClean : K := wide.eval (Environment.fromArray row fun _ _ ↦ #[])
def widePoly : K := (exprToCMv 3 wide).eval rowFn
#guard widePoly = wideClean

end Probe
\end{leancode}
\src{.claude/reports/blueprint-review/probes/boundary-adaptor/PolyBridge.lean}

\paragraph{Command} (from the repository root).
\begin{shellcode}
flock .claude/reports/blueprint-review/logs/lean.lock \
  lake env lean .claude/reports/blueprint-review/probes/boundary-adaptor/PolyBridge.lean
\end{shellcode}

\paragraph{Recorded output} (\file{PolyBridge.out}).
\begin{srccode}[numbers=left,numbersep=3pt]
def CPoly.toCMvPolynomial.{u_1} : {n : ℕ} →
  {R : Type u_1} → [inst : CommSemiring R] → MvPolynomial (Fin n) R → CMvPolynomial n R :=
fun {n} {R} [CommSemiring R] p =>
  match p with
  | AddMonoidAlgebra.ofCoeff { support := s, toFun := f, mem_support_toFun := mem_support_toFun } =>
    let unlawful := Unlawful.ofList (List.map (fun m => (CMvMonomial.ofFinsupp m, f m)) s.toList);
    ⟨unlawful, ⋯⟩
def CPoly.fromCMvPolynomial.{u_1} : {n : ℕ} →
  {R : Type u_1} → [inst : CommSemiring R] → CMvPolynomial n R → MvPolynomial (Fin n) R :=
fun {n} {R} [CommSemiring R] p =>
  let support := List.map CMvMonomial.toFinsupp (Lawful.monomials p);
  let toFun := fun f => p[CMvMonomial.ofFinsupp f]?.getD 0;
  have mem_support_fun := ⋯;
  AddMonoidAlgebra.ofCoeff { support := support.toFinset, toFun := toFun, mem_support_toFun := ⋯ }
exit=0
\end{srccode}

\subsection{\texorpdfstring{\protect\code{PolyBridgeNew}}{PolyBridgeNew}}\label{sec:gen-boundary-probe-polybridgenew}

\paragraph{What it establishes.} The same as \bndc{PolyBridge}
(\cref{sec:gen-boundary-probe-polybridge}) at the new pins, with \bndc{K.ofBits n} for the row's
numerals; every guard passes. \bndc{toCMvPolynomial} is still a \bndc{noncomputable def} at the new
CompPoly pin \bndc{572f9973} (\bndc{MvPolyEquiv/Core.lean:41}, read with \bndc{git show}).

\paragraph{Revision.} Run on 2026-09-30, after the rebuild, against the checkout \code{docs/protocol-blueprint-review} at \code{b022e4f}, whose Lean sources are \code{main} at \code{144c5aa}, with the new pins (CompPoly \code{572f9973}, Clean \code{42fe4b26}, ArkLib \code{7653a901}). Exit 0.

\paragraph{Source.}
\begin{leancode}[numbers=left,numbersep=3pt]
/-
  Probe (boundary-adaptor), NEW-PIN VARIANT (CompPoly 572f9973, K.ofBits in place of numerals): from a Clean expression to the polynomial an `M3Instance` holds.
  The blueprint's Layer 2 produces Mathlib's `MvPolynomial`; the instance holds CompPoly's
  `CMvPolynomial`. CompPoly's conversion from the first to the second is noncomputable at the
  pin; a direct, computable translation exists. Scratch work for the blueprint review.
  Run from the repository root:
    flock .claude/reports/blueprint-review/logs/lean.lock \
      lake env lean .claude/reports/blueprint-review/probes/boundary-adaptor/PolyBridgeNew.lean
-/
import LeanerVM
import CompPoly.Multivariate.MvPolyEquiv.Eval
import CompPoly.Multivariate.Operations

open LeanerVM.Parameters CompPoly CPoly

namespace Probe

-- The two directions of CompPoly's bridge at the pin.
#print CPoly.toCMvPolynomial
#print CPoly.fromCMvPolynomial

/-- A direct translation of a Clean expression over `K` to a computable polynomial in the first
`n` row variables; a variable past the width reads `0`, as Clean's `Environment.fromArray`
reads a missing cell. -/
def exprToCMv (n : ℕ) : Expression K → CMvPolynomial n K
  | .var v => if h : v.index < n then CMvPolynomial.X ⟨v.index, h⟩ else 0
  | .const c => CMvPolynomial.C c
  | .add a b => exprToCMv n a + exprToCMv n b
  | .mul a b => exprToCMv n a * exprToCMv n b

/-- The first `JUMP` residual, `b + v_cond · w`, on a row `(v_cond, w, b)`. -/
def residual : Expression K := .add (.var ⟨2⟩) (.mul (.var ⟨0⟩) (.var ⟨1⟩))

/-- A row. -/
def row : Array K := #[K.ofBits 3, K.ofBits 5, K.ofBits 7]

/-- The row as the assignment of the three variables. -/
def rowFn : Fin 3 → K := fun i ↦ row[i.val]?.getD 0

/-- Clean's value of the residual on the row. -/
def cleanValue : K := residual.eval (Environment.fromArray row fun _ _ ↦ #[])

/-- The polynomial's value on the row. -/
def polyValue : K := (exprToCMv 3 residual).eval rowFn

-- The translation runs, agrees with Clean's evaluation on the row, and has the degree the
-- expression shows.
#guard polyValue = cleanValue
#guard (exprToCMv 3 residual).totalDegree = 2

-- A variable past the width is read as zero on both sides.
def wide : Expression K := .add (.var ⟨5⟩) (.var ⟨0⟩)
def wideClean : K := wide.eval (Environment.fromArray row fun _ _ ↦ #[])
def widePoly : K := (exprToCMv 3 wide).eval rowFn
#guard widePoly = wideClean

end Probe
\end{leancode}
\src{.claude/reports/blueprint-review/probes/boundary-adaptor/PolyBridgeNew.lean}

\paragraph{Command} (from the repository root).
\begin{shellcode}
flock .claude/reports/blueprint-review/logs/lean.lock \
  lake env lean .claude/reports/blueprint-review/probes/boundary-adaptor/PolyBridgeNew.lean
\end{shellcode}

\paragraph{Recorded output} (\file{PolyBridgeNew.out}).
\begin{srccode}[numbers=left,numbersep=3pt]
def CPoly.toCMvPolynomial.{u_1} : {n : ℕ} →
  {R : Type u_1} → [inst : CommSemiring R] → MvPolynomial (Fin n) R → CMvPolynomial n R :=
fun {n} {R} [CommSemiring R] p =>
  match p with
  | AddMonoidAlgebra.ofCoeff { support := s, toFun := f, mem_support_toFun := mem_support_toFun } =>
    let unlawful := Unlawful.ofList (List.map (fun m => (CMvMonomial.ofFinsupp m, f m)) s.toList);
    ⟨unlawful, ⋯⟩
def CPoly.fromCMvPolynomial.{u_1} : {n : ℕ} →
  {R : Type u_1} → [inst : CommSemiring R] → CMvPolynomial n R → MvPolynomial (Fin n) R :=
fun {n} {R} [CommSemiring R] p =>
  let support := List.map CMvMonomial.toFinsupp (Lawful.monomials p);
  let toFun := fun f => p[CMvMonomial.ofFinsupp f]?.getD 0;
  have mem_support_fun := ⋯;
  AddMonoidAlgebra.ofCoeff { support := support.toFinset, toFun := toFun, mem_support_toFun := ⋯ }
exit=0
\end{srccode}

